\documentclass[dvisvgm,tikz,border=3pt]{standalone}
\usepackage{amsmath,amssymb}
\usepackage{pgfplots}
\usepackage{CJKutf8}
\pgfplotsset{compat=1.18}
\usetikzlibrary{arrows.meta,calc,positioning}

\definecolor{themebg}{HTML}{30362d}
\definecolor{themefg}{HTML}{edf4e8}
\definecolor{themegray}{HTML}{a4af9d}
\definecolor{themecurve}{HTML}{7eb6ff}
\definecolor{themealert}{HTML}{ff7b7b}
\definecolor{themeamber}{HTML}{f5b942}

\pgfdeclarelayer{background}
\pgfdeclarelayer{foreground}
\pgfsetlayers{background,main,foreground}

\begin{document}
\begin{CJK*}{UTF8}{gbsn}
\pagecolor{themebg}
\begin{tikzpicture}[color=themefg, text=themefg]

\begin{axis}[
    width=14.0cm,
    height=8.0cm,
    axis lines=middle,
    axis line style={color=themefg, -{Stealth}},
    tick style={color=themefg},
    tick label style={color=themefg, font=\scriptsize},
    every axis label/.style={color=themefg},
    xmin=-7.2, xmax=7.5,
    ymin=-2.0, ymax=2.8,
    xtick={-6.28, -3.14, 0, 3.14, 6.28},
    xticklabels={$-2\pi$, $-\pi$, $0$, $\pi$, $2\pi$},
    ytick={-1, 0, 1},
    yticklabels={$-1$, $0$, $1$},
    clip=false,
    xlabel={$x$},
    ylabel={$y$},
    title style={font=\footnotesize\bfseries, at={(0.5,1.06)}, anchor=south},
    title={Fourier 级数与 Dirichlet 收敛：周期延拓后在跳跃间断点处收敛于左右极限中点},
]

\begin{pgfonlayer}{background}
    % 主值区间 [-\pi, \pi] 极微弱底色
    \fill[themecurve!15!themebg] (axis cs:-3.14,-1.5) rectangle (axis cs:3.14,1.5);
    \draw[themecurve!70!themebg, dashed, line width=1.0pt] (axis cs:-3.14,-1.5) -- (axis cs:-3.14,1.5);
    \draw[themecurve!70!themebg, dashed, line width=1.0pt] (axis cs:3.14,-1.5) -- (axis cs:3.14,1.5);
    \node[text=themecurve, font=\scriptsize\bfseries] at (axis cs:0, 1.35) {主值区间 $[-\pi,\pi]$};
\end{pgfonlayer}

% 周期延拓波形 (方波)
\addplot[themecurve, line width=2.0pt, domain=-6.28:-3.14] {1};
\addplot[themecurve, line width=2.0pt, domain=-3.14:0] {-1};
\addplot[themecurve, line width=2.0pt, domain=0:3.14] {1};
\addplot[themecurve, line width=2.0pt, domain=3.14:6.28] {-1};

% 空心断点
\addplot[only marks, mark=o, mark size=2.5pt, line width=1.2pt, themecurve, fill=themebg] coordinates {
    (-6.28,1) (-3.14,1) (-3.14,-1) (0,-1) (0,1) (3.14,1) (3.14,-1) (6.28,-1)
};

\begin{pgfonlayer}{foreground}
    % Dirichlet 收敛中点 (实心高亮红点)
    \addplot[only marks, mark=*, mark size=3pt, fill=themealert, draw=themefg] coordinates {
        (-6.28,0) (-3.14,0) (0,0) (3.14,0) (6.28,0)
    };

    % 中点标注
    \node[text=themealert, font=\scriptsize\bfseries, anchor=north] at (axis cs:0,-0.2) {
        $S(0)=\frac{f(0^-)+f(0^+)}{2}=0$
    };
    \node[text=themealert, font=\scriptsize\bfseries, anchor=south west] at (axis cs:3.25,0.2) {
        $S(\pi)=0$
    };

    % 定理结论框放置在顶部纯净空间 (完全避让任何曲线)
    \node[text=themefg, font=\scriptsize, anchor=north west] at (axis cs:-7.2, 2.75) {
        \begin{tabular}{l}
        \textbf{Dirichlet 收敛定理}：\\
        1. 连续点：$S(x)=f(x)$\\
        2. 跳跃点：$S(x_0)=\frac{f(x_0^-)+f(x_0^+)}{2}$\\
        3. 周期端点：$S(\pm\pi)=\frac{f(\pi^-)+f(-\pi^+)}{2}$
        \end{tabular}
    };
\end{pgfonlayer}

\end{axis}
\end{tikzpicture}
\end{CJK*}
\end{document}
